\pagebreak

\chapter{Compiler Versions}
        
        The current version is 11.6.

        The principal enhancement from version 3.n to version 4.0 was the
        introduction of an extra scope based on source files. The {\bf include}
        statement has now been removed, along with {\bf globalvars} and {\bf
        privatevars}; {\bf module} and {\bf source} statements have been added in
        their place. {\bf nomodule} and {\bf nosource} statements have been added
        to ensure that files are included using the correct mechanism. The compiler
        is not backward compatible in this respect though usually changes required
        to existing sources are quite small and local.

        Version 4.1 had a new optional clause, the {\bf unless} construct, appended
        to most target control structures to allow premature termination when an
        exception expression is true. This also includes an optional code block
        to be executed when the exception occurs.
        
        Version 4.2 contains list (\autorefph{sec:sseclists}) and map
        (\autorefph{sec:ssecmaps}) types.
        
        Version 5.0 introduces variable {\bf attributes} which replace the
        arbitrary string of arguments previously required by inbuilt procedures
        {\bf input()}, {\bf externout()} and {\bf clock()}. There is a new mode
        {\bf output} and procedure {\bf externout()} has been deleted. Procedure
        {\bf attributes()} allows attributes to be applied to variables and
        function {\bf attributes()} provides a means of examining attributes.
        Variable attributes are used for -
        \blist
        \item   applying constraints to generated code elements, particularly
                input and output buffers and pads
        \item   controlling some aspects of code generation such as queue buffer
                depths and code optimisation
        \blend
        
        There are now two memory modes, {\bf cmemory} for combinatorial output
        memory and {\bf rmemory} for registered output memory, where previously
        there was a single mode ({\bf memory}) used by both types of memory. The
        inbuilt procedures {\bf cmemory()} and {\bf rmemory()} are still used as before
        to create combinatorial output and registered output memories respectively
        
        Version 5.0 also now takes note of the order of a subscript range. Previous
        versions sorted the range limits to smaller..larger but version 5.0
        uses the order specified, creating a descending order if required.
        
        Inbuilt procedures {\bf externout()} and {\bf buffersize()} have been
        deleted, inbuilt procedure {\bf fatal()} has been renamed {\bf exit()} and
        an inbuilt procedure {\bf texit()} (same as trace() followed by exit()) has
        been added.
        
        Variables of type "type" and type "str" can now be cross-assigned or
        initialised (no cast is needed). Variables of type "type" cannot otherwise
        be used where type "str" is accepted in which case the type can be
        converted by inbuilt function {\bf typestring()}, but this restriction
        may be lifted in the future.
        
        Version 6.0 requires Java 1.6 where previous versions used Java 1.5.
        
        Version 6.1 will handle Xilinx Coolrunner and XC95 CPLDs with a subset
        of modes and operations. Queue, cmemory, rmemory and selectvalue modes
        cannot be used. Target multiplication, division, remainder and
        variable shift operations are not implemented.
        
        Version 6.2 has undegone an internal code reorganisation and has
        a few obscure but important bug fixes.
        
        Version 6.3 has deprecated the following procedures -

        \begin{tabular}{| l | l |}
        \hline
        deprecated & replacement \\
        \hline
        rram()      & rmemory()    \\  
        rramread()  & rmemoryread()  \\
        rramwrite() & rmemorywrite() \\
        cram()      & cmemory()     \\ 
        cramwrite() & cmemorywrite() \\
        \hline
        \end{tabular}

        and the following functions -

        \begin{tabular}{| l | l |}
        \hline
        deprecated & replacement \\
        \hline
        rramout()   & rmemoryout()   \\
        cramread()  & cmemoryread()  \\
        \hline
        \end{tabular}

        This has been done to use identifiers more appropriate to the variable mode
        and at the same time to reverse the address type and data type arguments
        to the variable creation procedures to match the order of address and
        data arguments to the associated write procedures. The deprecated
        procedures and functions remain in the implementation.
        
        Version 6.4 supports Virtex-5, including gigabit transceivers.
        
        Version 6.5 contains a branch statement, the target equivalent of
        a switch statement. It also has new inbuilt functions timestring() and
        elapsedtime().
        
        Version 6.6 now allows pin location attributes to be any type as long
        as they only contain "str" primitives. Previously these were restricted
        explicitly to types "str" and "[]str" but now may be nested arrays and
        structs as long as all primitives are "str". This version also contains
        the added procedures and functions -

        \begin{tabular}{| l | l |}
        \hline
        split()     & split a string into substrings            \\
        getfield()  & get the nth field from a struct           \\
        numfields() & get the number of fields in a struct      \\
        isstruct()  & determine if a variable type is a struct  \\
        shell()     & execute a command in a shell              \\
        exec()      & execute a command                         \\
        env()       & return a map of environment variables     \\
        \hline
        \end{tabular}
        
        Version 6.7 now implements target fixed-point types and arithmetic
        operations. This version also changed floating point implementation from
        using a struct for representation to a proper type with inbuilt
        functions to extract the sign, biased exponent and mantissa. Functions
        added are -
        
        \begin{tabular}{| l | l |}
        \hline
        Offset()                & the offset of a fixed point value\\
        BiasedExponent()        & the biased exponent of a floating point value\\
        BiasedExponentWidth()   & the biased exponent width of a floating point value\\
        FloatIsNegative()       & the sign of a floating point value\\
        Mantissa()              & the mantissa of a floating point value\\
        MantissaWidth()         & the mantissa width of a floating point value\\
        \hline
        \end{tabular}
        
        Version 6.8 included the iswritten() inbuilt function.

        Version 7.0 has conditional compilation. It also has had the round()
        inbuilt function extended to provide rounding from fixed or ufixed
        to int or uint using round-to-even.
        
        Version 7.1 has substantially improved code optimisation. It also extends
        the 'continuous' attribute to queue mode variables.
        
        Version 7.2 has a new target construct, 'sync when () {}'.
        
        Version 7.3 allows target mode values in a compound value.
        
        Version 7.4 has improved implementation of combinatorial output
        memory. This is more efficient in some cases and also gives some
        more choices in the smaller queue buffer depths than previously.

        Version 7.5 has the following new features -
        \blist
        \item   static assignments \verb'+<-', \verb'-<-' and \verb'*<-' to
                add to, subtract from or multiply the LHS variable
        \item   immediate floating point values can be added to, subtracted
                from or assigned to target fixed and ufixed types
        \item   static variables can have an "ALU" attribute which will use
                a single adder/subtracter with an input selector to combine
                additions, subtractions, increments, decrements and
                assignments
        \item   an option to generate nested logic gates rather than expressions
                in LUTs for many low-level constructs to change the way
                the vendor place-and-route software optimises the logic
        \blend
        This version also adds a flip-flop to the output of every SRL16
        shift register to improve timing.

        Version 7.6 has undergone some improvement in netlist optimisation. It
        also now handles the Virtex-5 DSP48E block multiplier and provides for
        specific instantiation of the DSP48E. In the case of Virtex-2
        multipliers or Virtex-5 DSP48Es where these do not use the associated
        output register and their outputs drive flip-flops (static registers)
        with identical clock enables and resets, then the flip-flops are delated
        and the element registers are used instead. A Virtex5 static variable
        with DSP attribute set will also use the same DSP48E block for ALU
        logic.

        Version 8.0 has changed the way lists and maps are used. Lists and maps
        are now accessed using inbuilt procedures and functions. Also in this
        version functions can have the pointer dereference unary operator
        \verb'->' operator appended where the function returns a pointer. In
        addition in that case the function call can occur on the left side of an
        assignment statement.

        Version 8.1 has a reworked parser that now allows a function call with a
        pointer dereference followed by increment or decrement operator to be a
        statement. It also allows subscripts and fields on function call
        returns.

        Version 8.2 now supports Virtex-5 and Virtex-6. It has been upgrgaded to
        use the DSP blocks in Virtex-5 and Virtex-6 under control of DSP
        attributes on selectvalue, static and queue mode variables. For all FPGAs
        it also allows addition and subtraction operations to be combined in a
        single ALU for selectvalue, static and queue mode variables when directed
        by an ALU attribute. Explicit support for IDELAY, ODELAY, IDDR, ODDR and
        DSP48E has been added to the library xc5v.3pl (CRC16, CRC32 and SYSMON
        were already included).

        Version 8.3 contains the waitpri() construct. This is primarily intended
        for arbitration of conflicting writes to queue variables but can also be
        used to arbitrate between memory accesses and other statements. Call to
        priority() to create a priority mode variable no longer needs to be
        given a type in which case it will be given type "[20]log" by default.
        The inbuilt function isread() was added. This function is true for one
        clock cycle when its queue argument is read.

        Version 8.4 allows queues of depth 0, i.e. unbuffered queues. The
        principal use of these is to synchronise two parallel execution threads,
        optionally passing a datum in the process. The use of these is of
        necessity somewhat restricted.
        
        Version 8.5, when called with no source file command line argument,
        pops up a GUI to open the file. The GUI has buttons to set some
        options and a button to start (parsing, immediate execution and
        then code generation). The output appears in a scrollable text window.
        A file can be repeatedly compiled and executed by a single button
        press. Version 8.5 will run on Linux (or any Unix), MAC OSX and
        Windows.
        
        Version 8.6 uses CLB logic for selectvalue mode variables by default,
        whereas previously 3-state buffers were used. An attribute is used
        to force 3-state implementation if desired. An additional attribute
        is now available to determine the output level when no inputs are
        selected. Inbuilt procedures and functions tselect() and tselecta()
        have been removed.
        
        Version 8.7 allows selectvalue mode variables to have their default
        values, when no input is selected, to be specified. The environment
        variable {\bf THREEPL\_CONTINUOUS} is no longer recognised.
        There is a new directive, "output directory", produced by the
        compiler. The optional report file is now placed in the designated
        output directory instead of the parent directory.
        
        Version 8.8 has added inbuilt functions memoryaddresstype() and
        memorydatatype().
        
        Version 8.9 now allows compound values to contain target expressions.
        Inbuilt procedure float() will now accept an integer initialisation
        argument.
        
        Version 8.10 dispenses with the unworkable conditional compilation
        scheme and substitutes conditional file inclusion via additional
        forms of 'module' and 'source' statements.
        
        Version 8.11 no longer supports automatic conversion of "null" queues
        to "log" values. A queue of type "null" can be explicitly cast
        to a value of type "log" where required.
        
        Version 8.12 removes some restrictions on the use of unbuffered
        queues.
        
        Version 8.13 is functionally identical to version 8.12. The version
        number has been advanced because the source code has been extensively
        modified to incorporate generic Java typing and eliminate casts
        where possible. Also the parser has been separated from the main
        program, which was not the case previously.
        
        Version 8.14 removes the restriction on cmemory and rmemory variables
        being passed as arguments to varags parameters.
        
        Version 9.0 has removed clock managers and high speed serial devices
        from within 3PL and implemented them instead as 3PL library modules.
        This is a significant change and is not backward compatible.

	Version 9.1 checks for input variables with no clock domain that occur
	in target conditional statement test expressions.
        
        Version 10.0 introduces a major change. Most constraints and attributes
        have been moved from the netlist constraints file (NCF) to the
        EDIF netlist file where they now appear as element properties. Remaining
        constraints (principally timing) are still written to the NCF but
        are also translated into a Vivado XDC TCL script file. This means that
        there are no changes required to the source file or to 3PL run options
        or directives when changing between ISE and Vivado. Where constraints
        are generated in the 3PL source, mostly in library files, there are now
        calls to generate both NCF and XDC constraints separately. Library
        procedures INPUT(), OUTPUT() and CLOCK() have been introduced to
        properly handle constraints. These in turn call inbuilt procedures input(),
        output() and clock().
        
        Version 10.1 has the 'ignoremodules' option added to queue attributes.
        
        Version 10.2 allows map indices or keys to be applied in brackets in the
        same way as array subscripts. It also allows subscripts to be used to
        access lists, though list insertion still requires the use of an inbuilt
        procedure.
        
        Version 10.3 has changed inbuilt procedures input(), output() and clock()
        so that where possible I/O attributes (properties) normally contained in
        a constraint file are now handled in 3PL software, not through Java code
        as a side effect of netlist generation. Most of the attributes that these
        procedures accept are arrays, one member per pin, as this allows attributes
        to differ from pin to pin even within one variable.
        
        Version 11.0 changed all occurrences of 'pipe' to 'queue' in 3PL and
        the libraries. It was felt that naming a variable mode 'pipe' and then
        explaining that it was actually a queue was confusing and unnecessary.
        Also inbuilt procedures scanf, fscanf, sscanf, printf, fprintf and
        sprintf have been added.
        
        Version 11.1 allows static variables to be initialised via an "init"
        attribute. Inbuilt function strtotype() has been added.
        
        Version 11.2 no longer has inbuilt modules FD and REGISTER. There
        are numerous other minor code changes which should be backward
        compatible with version 11.1.
        
        Version 11.3 now allows many inbuilt modules and procedures and a few
        inbuilt functions to be called with parameter=argument pairs as well
        as matching arguments to parameters solely by position. Input mode and
        output mode variables no longer have clock domains. Period attributes
        have been removed, the reciprocal of frequency being used instead.
        
        Version 11.4 provides conditional compilation and execution based on
        time of last modification of all included source files, the compiler
        build time and the output netlist file. It also provides a
        3PL post-processing section where 3PL code can be executed after
        netlist generation is complete. This code can invoke place-and-route
        tools and any other post-processing operations which can again
        be conditional on the time of last modification of related files.

        Version 11.5 has had a minor change to command line options (-v
        no longer available). It has had a major restructuring of intermediate
        code signal structures. The report file format has been restructured
        and improved and the optional intermediate code list is now in
        a separate file.
        
        Version 11.6 has a change to the {\bf unless} control structure. It now
        requires an exception  code block in braces even if empty. The call to
        inbuilt procedure shell() has had its output argument order changed.
        
        Version 11.7 has classes but is otherwise backward compatible with
        version 11.6.
        
